\subsection{Features, predictors, and policies}
\label{sec:method}

\paragraph{Operational features.}
At prefix $t$, every learned head receives the same nine logit-free features.
We count searches used and distinct retained retrieval units. Question-term
coverage is the fraction of lowercased alphanumeric question tokens that
occur in $E_t$. New-evidence mass counts units added by the latest retrieval
after exact-string deduplication. Answer change compares normalized answers
after lowercasing and removal of punctuation and articles. Answer length is
the number of whitespace tokens. Grounding is the fraction of normalized
answer tokens found in $E_t$, with value zero for an empty answer or empty
evidence. A two-prefix stability indicator marks consecutive normalized-answer
agreement. The final feature counts distinct normalized answers in
$a_{0:t}$. These definitions make empty-prefix behavior deterministic and
keep the interface identical across backbones.

\paragraph{Predictors.}
Every prefix from a full-budget logged trajectory supplies graded FBG, binary
future-improvement, and answer-error labels. Where annotated support exists,
it also supplies an evidence-coverage label. Our FBG head is a
gradient-boosted regressor with 80 trees of maximum depth 2. Our error head is
logistic regression with \texttt{max\_iter=1000}. We fit the binary
future-improvement and gold-support classifiers with the same chain-level
cross-fitting discipline. Token F1 provides the two future-gain targets,
while EM remains our primary policy metric. The base LLM stays frozen and each
small predictor trains on CPU.

\paragraph{FBG policies.}
The FBG-only policy continues while $\fbgpred_t\ge\theta$ and stops at the
first crossing $\fbgpred_t<\theta$ or at $B$. The guarded FBG policy also
requires a latched first-repetition event. If the answer repeats at $t=1$,
changes at $t=2$, and the FBG score crosses at $t=3$, the policy stops at
$t=3$ because it retains the earlier repetition event. We use this explicit
name and timeline because the rule differs from requiring current
consecutive-answer stability. Prefix answers are already generated for the
features, so the guard adds no call under our answer-at-every-prefix schedule.

Algorithm~\ref{alg:voi} makes the latched score crossing, optional repetition
guard, and post-stop certificate explicit. The alternative policies below
replace the FBG score and crossing rule while retaining the same interaction
loop and budget cap.

\begin{algorithm}[t]
\DontPrintSemicolon
\SetKwInOut{Input}{input}
\SetKwInOut{Output}{output}
\caption{FBG-only and repeat-guarded retrieval control with an optional
post-stop certificate.}
\label{alg:voi}
\Input{goal $g$, frozen LLM, scores $(\fbgpred,\hat r)$, threshold $\theta$,
rule $\rho\in\{\mathrm{FBG},\mathrm{guarded}\}$, certificate cutoff
$q_\alpha$, budget $B$}
\Output{stopped answer $a_t$ and certificate flag}
$E_0\gets\varnothing,\quad \ell_{\mathrm{score}},\ell_{\mathrm{rep}}\gets\mathrm{false}$\;
\For{$t\gets0$ \KwTo $B$}{
  $a_t\gets\mathrm{LLMAnswer}(g,E_t)$\;
  $\ell_{\mathrm{score}}\gets \ell_{\mathrm{score}}\lor
  [\fbgpred(g,E_t,B-t,a_{0:t})<\theta]$\;
  $\ell_{\mathrm{rep}}\gets \ell_{\mathrm{rep}}\lor
  [t>0\land\mathrm{norm}(a_t)=\mathrm{norm}(a_{t-1})]$\;
  \If{$t=B\lor[\ell_{\mathrm{score}}\land
  (\rho=\mathrm{FBG}\lor \ell_{\mathrm{rep}})]$}{\textbf{break}}
  $q_t\gets$ LLM-proposed query given $(g,E_t)$\;
  $E_{t+1}\gets\operatorname{dedup}(E_t\cup\mathrm{search}(q_t))$\;
}
\Return{$\bigl(a_t,\mathbf{1}[q_\alpha\ne\bot\land
\hat r(g,E_t,B-t,a_{0:t})\le q_\alpha]\bigr)$}
\end{algorithm}

\paragraph{Alternative policies and selection.}
The binary future-improvement policy stops at the first nonterminal prefix
whose predicted probability is below its threshold. The error-head gate stops
when predicted current-answer error is at most its threshold, and the
gold-support policy stops when predicted evidence coverage is at least its
threshold. Every rule may stop at $t=0$ and defaults to $B$ if no threshold is
crossed. These directions matter because the heads estimate different events.

For every learned policy, we choose the cheapest cross-fitted development
point whose EM lies within one point of full-budget development EM. If no
point enters that band, we choose the highest-EM point and break ties by cost.
The fixed comparator is selected from the same development questions as the
smallest budget attaining maximum development EM. The learned rule and fixed
budget therefore receive comparable access to development outcomes, but the
objectives are intentionally asymmetric. The one-EM band asks how cheaply a
learned rule can remain near full-search quality, so it favors learned rules
on cost by permitting a small development loss. The fixed rule is more
quality-conservative because it first maximizes development EM and only then
minimizes cost. We retain $B^\star$ as a deterministic deployment reference,
then remove dependence on this asymmetry through the stronger matched-cost
hull comparison against every eligible simple operating point. Within each
signal family, all primary cells use the same crossing direction, declared
candidate grid, cross-fitting procedure, and selection code. Grids differ
across signal families only because their score scales differ. Our optional
certificate is applied only after the selected policy stops and never changes
$T$.
